% ===================================================================
%  第 7 节 魔群月光（Monstrous moonshine）
%  原文：TASI10.tex 第 2136--2179 行（PDF p.58--59）
% ===================================================================
\section{魔群月光}
\label{sec:physics_monstrous}

\begin{keytakeaways}
  \item 物理构造的起点是杂化弦紧化到 $1+1$ 维：内部 CFT 的中心荷为 $(24,12)$，由魔群月光模 $V^\natural$ 与 Conway 模 $V_{}^{s\natural}$ 的乘积给出；再紧化空间方向到圆，便得到 $0+1$ 维理论。
  \item 对每个 $g\in\MM$ 作 CHL 型轨道流形化，其第一量子化 BPS 态构成广义 Kac–Moody 代数 $\mathfrak{m_g}$ 的模：$g=e$ 时是 Borcherds 的魔代数，$g\neq e$ 时是 Carnahan 的 twisted 魔代数。
  \item BPS 指标 $Z_{(g,h)}(T,U)$ 可以写成 McKay–Thompson 级数的乘积（\eq{BPSindexmonster}），从而把时空 BPS 态计数与月光函数联系起来。
  \item BPS 指标在自对偶群下的不变性直接给出 $T_g$ 的 $\Gamma_g$ 不变模性；对半径插值的分析又给出 Hauptmodul 性质。
  \item 这条思路为魔群月光的模性与亏格零性质提供了物理解释，是“用弦论实现月光”的范式。
\end{keytakeaways}

本节我们简要描述近期的一个弦论构造，它由 \cite{paquette2016monstrous,Paquette:2017xui} 发展而来，
涉及魔群月光模 $V^\natural$ 以及 \S\ref{sec:mm} 讨论的相伴的扭模（twisted module）$V^\natural_g$。
尽管魔群月光令人瞩目的亏格零性质已在数学上得到证明，其物理解释长期难以捉摸。
近期工作 \cite{paquette2016monstrous,Paquette:2017xui} 通过考虑某些杂化弦理论紧化到 $0+1$ 维的情形，
为这一重要性质提供了线索。

% [Tutor Note: 这里的 $V^\natural_g$ 不是某个独立构造的代数，而是 $V^\natural$ 在 $g$ 作用下的扭模：其分次特征标正是 McKay–Thompson 级数 $T_g$。理解本节的关键是把“群元 $g$ 标注的模块”与“时空 BPS 态空间”对应起来，而不是把 $g$ 当成对称性作用在未扭空间上。]

在 \cite{paquette2016monstrous,Paquette:2017xui} 中，出发点是紧化到 $1+1$ 维，
其内部 CFT 的中心荷为 $(24,12)$，由 $V^\natural\times V_{}^{s\natural}$ 给出。
第二个因子是 \S\ref{sec:conway_moonshine} 引入的 Conway 模，它主要扮演旁观者的角色，
其主要功能是为世界面理论提供右行 $\mathcal{N}=1$ 超对称。
该模型在二维时空中具有 ${\cal N}=(0,24)$ 超对称。
我们进一步把剩余的空间方向紧化在圆上。
这一额外紧化对于计算下文描述的“BPS 指标”以及构造更一般的模型至关重要。

% [Tutor Note: 中心荷 $(24,12)$ 应读作左行 $c=24$、右行 $c=12$：$V^\natural$ 贡献左行 $24$，Conway 模 $V_{}^{s\natural}$ 贡献右行 $12$。右行 $\mathcal{N}=1$ 世界面超对称由 Conway 模提供，这正是后文能用 $(-1)^F$ 定义时空 BPS 指标的前提。]

更一般地，对每个 $g\in\MM$，通过按一个对称性对理论作轨道流形化来构造 CHL 型模型：
该对称性在 $V^\natural$ 上按 $g$ 作用，在圆上则按一个移位作用，其阶由 $g$ 在 $V^\natural$ 上作用的某些性质决定
（在最简单的情形中等于 $g$ 的阶）。
为避免涉及引力反常的微妙之处，该理论应在零弦耦合极限下考虑。

% [Tutor Note: CHL 型模型指 Chaudhuri–Hockney–Lykken 型紧化：在圆上同时作平移与对称性投影，使左右行不对称。圆上移位的阶必须与 $g$ 在 $V^\natural$ 上的作用阶匹配，否则扭结扇区的基态能量不为零、模不变性被破坏；这是构造能自洽的关键约束。]

\cite{paquette2016monstrous} 表明，上述模型的（第一量子化）BPS 态构成广义 Kac–Moody 代数 $\mathfrak{m_g}$ 的模：
当 $g= e$ 时它就是 Borcherds 构造的魔代数，当 $g\neq e$ 时则是 Carnahan 构造的“twisted”魔代数
\cite{Carnahan2008,MR2904095,carnahan2012generalized}。
广义魔群月光的讨论另见 \S\ref{sec:mm}。
我们关心这些理论中计数时空 BPS 态的（twined）BPS 指标。
记相关 Hilbert 空间为 $\mathcal{H}^g_{\text{BPS}}$，我们定义
\begin{equation}
Z_{(g,h)}(T,U):=\tr_{\mathcal{H}^g_{\text{BPS}}}\left((-1)^F \, h \, e^{2\pi iTW}e^{2\pi iUM}\right)~,
\end{equation}
其中 $g,h$ 是魔群中相互对易的元素。

% [Tutor Note: 指标 $Z_{(g,h)}$ 的两个下标含义不同：第一个下标 $g$ 选定扭结扇区，即迹在 $\mathcal{H}^g_{\text{BPS}}$ 上取；第二个下标 $h$ 是插入的群元。只有 $g$ 与 $h$ 对易，$h$ 才保持 $g$-扭结扇区，迹才有定义；这就是原文强调 commuting elements 的原因。]

上式中 $T,U\in \HH$ 分别是绕数（winding）与动量（momentum）量子数 $W,M$ 的复化化学势，
其虚部按逆温度与半径缩放为
\be
\im T \sim {\beta R}~,~~ \im U \sim {\beta \over R}~.
\ee

% [Tutor Note: 指数中 $T$ 配绕数 $W$、$U$ 配动量 $M$，不要写反。绕数能量正比于半径 $R$，动量能量正比于 $1/R$，所以 $\im T\sim\beta R$、$\im U\sim\beta/R$；这里 $T,U$ 是复化化学势，实部给出相位，虚部给出 Boltzmann 型压制。]

利用 $\mathcal{H}^g_{\text{BPS}}$ 的 Fock 空间构造（允许任意数目的自由弦），
以及把无穷乘积与无穷和联系起来的 twisted 分母公式——它推广了 Koike–Norton–Zagier 公式 \eq{eq:replicability_J}——可以得到
\begin{equation}\label{BPSindexmonster}
\begin{split}
&Z_{(e,g)}(T,U)=\left(T_{(e,g)}(T)-T_{(e,g)}(U)\right)^{24} \\
&Z_{(g,e)}(T,U)=\left(T_{(e,g)}(T)-T_{(g,e)}(U)\right)^{24}~.
\end{split}
\end{equation}

% [Tutor Note: 原文如此：第二行右侧为 $T_{(g,e)}(U)$，与第一行 $T_{(e,g)}(U)$ 的下标不对称；可能另有约定（扭模的 twining 方向），也可能是笔误。译文照录，读者引用时请核对原文。阅读时注意 $T_{(e,g)}=T_g$ 才是 McKay–Thompson 级数，不要把 $T_{(g,e)}$ 也当成 $T_g$。]

% [Tutor Note: “twisted 分母公式”指 Borcherds 的 twisted denominator identity（广义 Kac–Moody 代数的 Weyl 分母公式）：它把扭结扇区上的无穷乘积展开为无穷和，其指数正是 $\mathfrak{m_g}$ 的根重数。Koike–Norton–Zagier 公式 \eq{eq:replicability_J} 是 $g=e$ 时的特例；作者略去了该恒等式的证明。]

注意 $T_{(e,g)}=T_g$ 是魔群月光的 McKay–Thompson 级数。

上述表达式给出魔群月光的模性与亏格零性质的物理解释。

% [Tutor Note: $Z_{(g,e)}$ 取 $1/24$ 次幂后与 $T_g$ 同阶：$Z_{(g,e)}(T,U)^{1/24}$ 在 $T\to i\infty$ 时的发散对应 $T_g$ 的极点项 $q^{-1}$。这个 $1/24$ 次幂来自 24 个横方向自由弦因子的乘积，指数极易漏掉；不要把 $1/24$ 误写成 $1/2$ 或 $24$。]

首先，\cite{paquette2016monstrous} 利用 \cite{Persson:2015jka} 发展的技巧表明，
自对偶群（即 T 对偶群的一个子群，它在同一个 CHL 模型之间作映射，但一般在不同的模值处）
涉及 $\SL_2(\RR)$ 的两个子群，它们在 \eq{BPSindexmonster} 中分别独立地作用于 $T$ 与 $U$。
特别地，作用于 $T$ 的 $\SL_2(\RR)$ 子群正是 $\Gamma_g$，即魔群月光中相应的亏格零群
（参见猜想 \ref{conj:conwa_norton}）。
因此，BPS 指标在自对偶群下的物理不变性立即给出月光函数 $T_{g}(T)$ 的 $\Gamma_g$ 不变模性。

% [Tutor Note: 自对偶群是 T 对偶群中保持同一个 CHL 模型（可以取不同模值）的子群；它能分解为分别作用在 $T$ 与 $U$ 上的两个 $\SL_2(\RR)$ 子群，这是二维环面紧化的结构使然。作用在 $T$ 上的那个子群恰为亏格零群 $\Gamma_g$，所以指标的模不变性“翻译”成 $T_g$ 的模性；注意这里不需要证明 $T_g$ 模性，它是物理不变性的推论。]

其次，回忆 $T_{(e,g)}$ 的 Hauptmodul 性质可以重新表述为：
它们只在 $\tau\rightarrow i\infty$ 及其在 $\Gamma_g$ 下的像处有单个极点。
用超对称指标的语言，这意味着对固定的 $U$，$Z_{(g,e)}(T,U)^{1/24}$ 及其 T 对偶像在 $T\rightarrow i\infty$ 时发散。
利用 T 对偶
\[
Z_{(e,g)}(T,U)=-Z_{(e,g)}\left(-\frac{1}{U},-\frac{1}{T}\right)
\]
以及 $T_{(g^ah^c,g^bh^d)}(\tau)$ 与 $T_{(g,h)}\left(\frac{a\tau+b}{c\tau+d}\right)$ 之间的广义月光关系
（参见 \S\ref{sec:mm}），这一性质就转化为如下陈述：
对所有互素整数 $a$ 与 $b$，当 $T$ 固定时 $Z_{(g^a,g^b)}(T,U)$ 在 $U\to i\infty$ 的极限下发散，
当且仅当该极限通过某个与 $g\in \MM$ 相联系的 CHL 模型的自对偶，
与 $T\to i\infty$（$U$ 固定）的极限相联系。

% [Tutor Note: 上式 T 对偶的整体负号极易漏掉：$Z_{(e,g)}(T,U)=-Z_{(e,g)}(-1/U,-1/T)$ 同时交换了 $T$ 与 $U$、各自取倒数，并带一个负号。注意 $(f(T)-f(U))^{24}$ 在 $T\leftrightarrow U$ 下因指数为偶数并不变号，所以这个负号是原文式子的独立成分，照录即可，不要自行“修正”为加号；若改成加号，后文用 T 对偶连接两个极限的论证将不成立。]

下面我们从物理考虑论证这一点为何成立。
让我们通过改变半径 $R$ 在两个极限之间插值，同时把逆温度 $\beta$ 固定得很大。
在低温极限下，固定 $U$、$T\to i\infty$ 处的发散由基态贡献 $e^{\beta R}$ 引起，它在 $R\to \infty$ 时爆炸。
小半径极限是收敛的，除非发生相变，使得在临界线的另一侧出现贡献为 $e^{\beta / R}$ 的新基态。
若发生这种情况，可以证明在临界线上的自对偶半径极限处存在 $SU(2)$ 对称性，
它特别包含一个 T 对偶变换。
换言之，临界线两侧由包含在 $\Gamma_g$ 中的一个自对偶相联系，这就证明了 Hauptmodul 性质。

% [Tutor Note: 作者在此跳过论证的一步是：自对偶半径 $R=\sqrt{\alpha'}$ 处，圆上的 $U(1)_L\times U(1)_R$ 会增强为 $SU(2)_L\times SU(2)_R$（本情形只剩一个 $SU(2)$），其 Weyl 群包含 $R\leftrightarrow 1/R$ 的 T 对偶。正是这个增强对称性保证相变两侧的基态由自对偶联系，从而把“单极点”翻译成 Hauptmodul 性质；完整论证见 \cite{paquette2016monstrous}。]

% [Tutor Note: “$0+1$ 维”即 0 个空间维加 1 个时间维（量子力学）；原文第一部分引言处写作 “$(1+0)$-dimensional”，两种写法指同一次紧化，不要误认为先紧化到两个不同维数。紧化链是：10 维杂化弦 → 2 维（内部 $(24,12)$ CFT）→ 圆紧化 → $0+1$ 维。]
